\section*{第 \ref{ch:g12:rlc-oscillations} 章——自由电振荡：RLC}

\begin{solution}{exo:g12:rlc-oscillations:1}
(a) $\sqrt{LC} = \qty{1.0e-4}{s}$: $T_0 = \qty{0.63}{ms}$,
$f_0 \approx \qty{1.6}{kHz}$. (b) $\sqrt{LC} = \qty{3.16e-7}{s}$:
$T_0 = \qty{2.0}{\micro s}$，$f_0 \approx \qty{0.50}{MHz}$——相差
$\sqrt{10^5} \approx 316$ 倍。
\end{solution}

\begin{solution}{exo:g12:rlc-oscillations:2}
$LC$: $(\unit{V.s/A})(\unit{C/V}) = \unit{s} \times \unit{C/A} =
\unit{s} \times \unit{s} = \unit{s^2}$，故 $\sqrt{LC}$ 以秒为单位。
$2\pi$ 是纯数——没有量纲需要改变。
\end{solution}

\begin{solution}{exo:g12:rlc-oscillations:3}
$Q_0 = CU_0 = \qty{2.8}{\micro C}$; $E = \tfrac12 CU_0^2 =
\tfrac12 \times \num{4.7e-7} \times 36 \approx \qty{8.5}{\micro J}$.
电荷以 $f_0 = 1/(2\pi\sqrt{LC})$ 的频率振荡，电流
与之相差四分之一周期，能量在极板与线圈之间来回转移，总
能量保持不变。
\end{solution}

\begin{solution}{exo:g12:rlc-oscillations:4}
$Q_0 = \qty{2.0}{\micro C}$; $f_0 = 1/T_0 = \qty{250}{Hz}$;
$q(T_0/2) = \qty{-2.0}{\micro C}$；$i(0) = 0$（余弦在 $t=0$ 处取极值）；
$I_0 = 2\pi \times \num{2.0e-6}/\num{4.0e-3} \approx \qty{3.1}{mA}$.
\end{solution}

\begin{solution}{exo:g12:rlc-oscillations:5}
能量全部在电容器中时，$q = \pm Q_0$，$i = 0$（$t = 0$、$T_0/2$、
$T_0$、\dots），相当于单摆的两个端点。能量全部在线圈中时，$q = 0$，
$i = \pm I_0$（$t = T_0/4$、$3T_0/4$、\dots）时，相当于单摆的
最低点，也是速度最快的点。
\end{solution}

\begin{solution}{exo:g12:rlc-oscillations:6}
$d^2q/dt^2 = -(2\pi/T_0)^2\,Q_0\cos(2\pi t/T_0) = -(2\pi/T_0)^2 q$,
于是 $L\,d^2q/dt^2 + q/C = q\,[\,1/C - L(2\pi/T_0)^2\,] = 0$ 对所有 $t$
成立，当且仅当 $(2\pi/T_0)^2 = 1/(LC)$，即 $T_0 = 2\pi\sqrt{LC}$。
\end{solution}

\begin{solution}{exo:g12:rlc-oscillations:7}
$L = \dfrac{1}{4\pi^2 f_0^2 C} = \dfrac{1}{39.5 \times
(\num{1.98e5})^2 \times \num{2.2e-10}} \approx \qty{2.9}{mH}$.
\end{solution}

\begin{solution}{exo:g12:rlc-oscillations:8}
$E = \tfrac12 CU_0^2 = \qty{5.0e-5}{J}$；$\tfrac12 LI_0^2 = E$ 得
$I_0 = \sqrt{2E/L} = \sqrt{\num{1.0e-2}} = \qty{0.10}{A}$，首次
出现于 $T_0/4 = \tfrac14 \times \qty{0.63}{ms} \approx
\qty{0.16}{ms}$.
\end{solution}

\begin{solution}{exo:g12:rlc-oscillations:9}
$E_C \propto \cos^2(2\pi t/T_0)$，而 $\cos^2$ 每
\emph{半}圈重复一次：每个周期内电容器被充满两次，故
频率为 $2f_0$。当 $t = T_0/8$ 时，相位为 $\pi/4$：
$\cos^2(\pi/4) = \tfrac12$，故 $E_C = E_L = E/2$。
\end{solution}

\begin{solution}{exo:g12:rlc-oscillations:10}
每个准周期的比值为 $3.0/4.0 = 0.75$，第六个极大值
为 $4.0 \times 0.75^6 \approx \qty{0.71}{V}$。能量 $\propto$ 振幅
的平方，因此每个周期约保留 $0.75^2 \approx 56\%$ 的能量。准周期
仍 $\approx T_0 = 2\pi\sqrt{LC}$：轻阻尼只使振幅减小，
而几乎不影响振动快慢。
\end{solution}

\begin{solution}{exo:g12:rlc-oscillations:11}
$T_0 = 2\pi\sqrt{0.50 \times \num{2.0e-5}} = 2\pi \times
\num{3.16e-3} \approx \qty{20}{ms}$。对应关系：$k = m/(LC) =
0.50/\num{1.0e-5} = \qty{5.0e4}{N/m}$；检验：
$2\pi\sqrt{m/k} = 2\pi\sqrt{\num{1.0e-5}}$——结果同样为 \qty{20}{ms}。
\end{solution}

\begin{solution}{exo:g12:rlc-oscillations:12}
$C = \qty{500}{pF}$: $f_0 = 1/(2\pi\sqrt{\num{1.0e-13}}) \approx
\qty{0.50}{MHz}$; $C = \qty{50}{pF}$: $f_0 \approx \qty{1.6}{MHz}$.
$L$ 固定时，$f_0 = (2\pi\sqrt{L})^{-1} C^{-1/2} \propto 1/\sqrt{C}$：
$C$ 增大十倍，只以 $\sqrt{10} \approx 3.2$ 倍的因子改变 $f_0$——与前面结果一致。
\end{solution}

\begin{solution}{exo:g12:rlc-oscillations:13}
$0.95^n = \tfrac12$: $n = \ln 2/\lvert\ln 0.95\rvert \approx 13.5$,
故约 $14$ 个周期。振幅 $\propto \sqrt{E}$：降至
初始值的 $1/\sqrt{2} \approx 0.71$ 倍。若频率为 \qty{1.0}{MHz}，
$T_0 = \qty{1.0}{\micro s}$，约为 \qty{14}{\micro s}。
\end{solution}

\begin{solution}{exo:g12:rlc-oscillations:14}
每周期损耗 $0.020 \times \qty{1.0e-9}{J} = \qty{2.0e-11}{J}$，
按每秒 \num{32768} 个周期计，$P = \num{32768} \times \num{2.0e-11}
\approx \qty{6.6e-7}{W}$——不足一微瓦，纽扣
电池可用数年。同相推动时，只有在正确相位施力才会\emph{增加}
能量（共振）；反相则会阻碍摆动。
\end{solution}

\begin{solution}{exo:g12:rlc-oscillations:15}
$C = \dfrac{1}{4\pi^2 f_0^2 L} = \dfrac{1}{39.5 \times
(\num{2.5e7})^2 \times \num{1.0e-6}} \approx \qty{41}{pF}$。电阻器
应使回路进入非周期状态：电荷
随后无过冲地趋于稳定——既无振荡，也无振铃。
\end{solution}

\begin{solution}{pb:g12:rlc-oscillations:1}
\textbf{1.} 回路定律：$u_L + u_C = 0$，其中 $u_C = q/C$，
$u_L = L\,di/dt$, $i = dq/dt$: $L\,d^2q/dt^2 + q/C = 0$.

\textbf{2.} $d^2q/dt^2 = -(2\pi/T_0)^2 q$，故方程化为
$q\,[\,1/C - L(2\pi/T_0)^2\,] = 0$，对所有 $t$ 成立，当且仅当
$T_0 = 2\pi\sqrt{LC}$.

\textbf{3.} $LC$：$(\unit{V.s/A})(\unit{C/V}) = \unit{s} \times
\unit{C/A} = \unit{s^2}$——$\sqrt{LC}$ 是一个时间。

\textbf{4.} $LC = \qty{6.6e-14}{s^2}$：$T_0 = 2\pi \times \num{2.57e-7}
\approx \qty{1.6}{\micro s}$，$f_0 \approx \qty{620}{kHz}$——位于
\qtyrange{530}{1700}{kHz} 之内。

\textbf{5.} $C = \dfrac{1}{4\pi^2 f_0^2 L} = \dfrac{1}{39.5 \times
\num{1.0e12} \times \num{2.0e-4}} \approx \qty{127}{pF}$。

\textbf{6.} 用同一公式，取 \qty{530}{kHz}：$C \approx \qty{451}{pF}$。

\textbf{7.} 取 \qty{1700}{kHz}：$C \approx \qty{44}{pF}$。

\textbf{8.} $C = \qty{480}{pF}$：$f_0 \approx \qty{514}{kHz}$；
$C = \qty{20}{pF}$：$f_0 \approx \qty{2.5}{MHz}$——覆盖整个波段，
两端都留有余量。

\textbf{9.} $f_0 = (2\pi\sqrt{L})^{-1}C^{-1/2} \propto 1/\sqrt{C}$；
要使 $f_0$ 加倍，须将 $C$ 除以 $4$。

\textbf{10.} 由于 $f_0 \propto 1/\sqrt{C}$，等量改变 $C$ 时，
$f_0$ 在 $C$ 较大时变化缓慢，在 $C$ 较小时变化迅速：于是
相邻电台之间的间隔只需很小的调谐旋钮转动量——它们挤在
高频（即 $C$ 较小）一端。

\textbf{11.} $Q_0 = CU_0 = \num{1.27e-10} \times 0.020 \approx
\qty{2.5e-12}{C}$；$E = \tfrac12 CU_0^2 \approx \qty{2.5e-14}{J}$。

\textbf{12.} $\tfrac12 LI_0^2 = E$：$I_0 = U_0\sqrt{C/L} =
0.020\sqrt{\num{6.35e-7}} \approx \qty{16}{\micro A}$。

\textbf{13.} $0.96^n = \tfrac12$：$n = \ln 2/\lvert\ln 0.96\rvert
\approx 17$ 个周期。

\textbf{14.} 取 \qty{1.00}{MHz}，$T_0 = \qty{1.0}{\micro s}$：振铃
约在 \qty{17}{\micro s} 内减半。单个振铃在微秒量级即消失；
要有音乐，就须让回路持续获得能量补充。

\textbf{15.} 电波每秒推动数百万次，只有当
其频率与 $f_0$ 匹配时，各次推动才会同步到达并使振荡逐渐增强——
这就是共振；失谐的电台推动时，同向与反向的次数相同，
相互抵消。在发射端，放大器每个周期补回耗散掉的能量：
由此得到一个受维持的振子。

\textbf{16.} $LC = \qty{2.54e-14}{s^2}$：$k = m/(LC) =
0.20/\num{2.54e-14} \approx \qty{7.9e12}{N/m}$。

\textbf{17.} 约为汽车弹簧的 $10^8$ 倍：没有任何台式
弹簧振子能达到兆赫兹——日常的劲度系数与
质量都慢得毫无希望。

\textbf{18.} $m = \dfrac{k}{(2\pi f_0)^2} = \dfrac{100}
{(\num{6.28e6})^2} \approx \qty{2.5e-12}{kg}$——仅几纳克，
相当于一粒尘埃。

\textbf{19.} 一块石英晶体：质轻如羽、刚度极高的薄片，
以其固有频率振鸣——即本学年前文提到的
计时器。

\textbf{20.} 取 $L = \qty{0.20}{mH}$，电容在 \qtyrange{20}{480}{pF} 之间变化时，回路
扫过约 \qtyrange{0.51}{2.5}{MHz}，覆盖整个波段；
$C \approx \qty{127}{pF}$ 时选中 \qty{1.00}{MHz}；若任其自由，振铃
约在 $\approx \qty{17}{\micro s}$ 内减半，但电台自身的电波，
与之同步共振，重新为回路供能——新闻播音依靠的
正是电波本身的能量。
\end{solution}
